\section{The Priest and the King}

\begin{quotex}
There is nothing that can be truly and well done or made except by the man in whom the marriage of the Sacerdotium and the Regnum has been consummated, nor can any peace be made except by those who have made their peace with themselves.
\flright{\textsc{Ananda Coomaraswamy}}
\end{quotex}


In 1942, \textbf{Ananda Coomaraswamy} (AKC) followed up his earlier review of \textbf{Julius Evola\footnote{\url{https://www.gornahoor.net/?p=4274}}'s Revolt Against the Modern World}. His short book Spiritual Authority and Temporal Power in the Indian Theory of Government simultaneously builds on \textbf{Rene Guenon}'s book of a similar title and exposes the misunderstanding underlying Evola's position.

Let us be clear. My purpose is not in the academic style which would state AKC's and Evola's positions for the purpose of comparison. Quite the contrary, it is necessary once and for all to show that Evola's position leads to confusion and hence must be rejected. The reason is unassailable since it is a question of a metaphysical principle, although historical examples may be illustrative.

Evola's errors arise from failing to understand the metaphysical functions of castes, confusing spiritual authority with temporal power, and seeing masculinity and femininity in absolute terms rather than relative.

\paragraph{Power and Authority}


In his discussion of the degeneration of castes (ch 35 of Revolt), Evola does not always clearly distinguish power and authority. He sees this degeneration beginning when the Kshatriya caste seized power from the Brahmins. However, that was never their proper relationship. The Brahmins are the custodians of spiritual authority and their role is not the exercise of power of temporal affairs. The Brahmins are related to the Kshatriya as knowledge to action, the mover to the moved, the interior to the exterior. The revolt of the Kshatriyas, therefore, does not consist in seizing power, but rather in the rejection of the spiritual authority that guides and is the principle of action.

\paragraph{Coomaraswamy's Critique}


AKC begins the book with a reference to Revolt and to Evola's misunderstanding of an event in the Brahamana. AKC again spoke highly of the chapter on \emph{Man and Woman}, although its effect is diminished by Evola's misunderstanding. The passage concerns the marriage formula that the groom says to the bride:

\begin{quotex}
I am That, thou art This, I am Sky, thou art Earth.
\end{quotex}


When the king chooses his Purohita (spiritual adviser), the Priest addresses the same formula to the king. To be clear, this means that the Priest is masculine in respect to the king, the feminine party in that relationship. When Evola was informed of this, he apparently dropped the reference to this passage in subsequent editions of Revolt without altering his claim. I will quote AKC in full, keeping in mind that Sacerdotium refers to the spiritual domain and Regnum to the temporal.

\begin{quotationx}
Evola's thesis, in his discussion of the Regnum, forces him to misinterpret this marriage formula . Had it not been for this, his admirable chapter ``Man and Woman'', applied to the true relationships of the Sacerdotium and the Regnum (approximately ``Church and State''), would have acquired a greater significance. As it is, Evola's argument for the superiority of the Regnum, the active principle, to the Sacerdotium, the contemplative principle, is a concession to that very ``modern world'' against which his polemic is directed.

His argument is as much a perversion of the Greek as it is of the Indian doctrine. In the Greek tradition the heroic kind of caste, alike in the soul and the community, ``that part of our soul which is endowed with bravery and courage, and which is the lover of victory,'' (Timaeus) is the best part of the mortal or animal soul, superior to the appetitive and inferior to the spiritual and immortal part that lays down the law. As such its seat is in the heart between the bowels and the head; it is the defender of the whole community; its function is to listen to the Voice Logos from the Acropolis, to serve and cooperate in battle with the sacred principle against the mob of the appetites (within us) or of moneyed men (in the city). The three parts of the soul (or body politic) thus evidently correspond in hierarchy to the brahmin, kshatriya, and vaishya, respectively the Sacerdotium, Regnum, and Commons of the Vedic tradition (in which the shudra is represented by the Asuras); and there can be no possible doubt of the superiority of the sacred to the royal character.

That the Spiritual Authority, Plato's \emph{leron}, is also the Ruler, Plato's \emph{archon}, just as the brahmin is ``both the brahmin and the Kshatriya,'' means indeed that the Supreme Power is a royal as well as a priestly power, but quite certainly does not mean that the Kshatriya considered apart from the brahmin is itself the supreme authority or anything more than its agent and servant.
\end{quotationx}


AKC goes on for 87 pages with densely packed allusions to Hindu mythology that make this point in increasing detail; further discussion will have to wait for another day.

\paragraph{Warrior and Ascetic}


In a telling passage in ``Man and Woman'', Evola expresses AKC's point:

\begin{quotex}
According to metaphysical symbols, the female becomes the ``bride'' which is also the ``power'' or instrumental generating force that receives the first principle of the immobile male's impulse and form.
\end{quotex}


Since the king and the Kshatriya are the power wielders, Evola inadvertently illustrates AKC's point that the king is the bride in relation to the priest. Shortly after that, Evola reverts back and repeats his error

\begin{quotex}
The mode of being that corresponds principally to man was already considered; and we said on the two principle forms closest to the value of ``being in oneself'': Action and Contemplation, the Warrior (or Hero) and the Ascetic are therefore the two fundamental types of pure virility.
\end{quotex}


Although both these types may represent virility, they are not coequal and independent. Only the Contemplative has the principle of his being in himself. The Warrior's first principle, as has already been shown, lies outside himself, and specifically in knowledge. That does not mean the path of the warrior is ineffective. For the Contemplative, there is the path of jnana yoga and for the Warrior, there is bhakti yoga. As Evola writes,

\begin{quotex}
It is a norm that obeys the same principle as the caste system, and it refers to the two cornerstones of \emph{dharma} and \emph{bhakti}, or \emph{fides}: one's own nature and active devotion.
\end{quotex}


This is precisely AKC's point about Evola's chapter on man and woman. The principle that relates man and woman is the same one as in the caste system, just as Evola says. Hence, the relationship between the Priest, or Brahmin, and the King, or Kshatriya, is analogous to that between man and woman.

We see again that the rejection of a metaphysical principle is not like the rejection of a geometric postulate that produces a new, yet logically consistent, geometry. Rather, it leads to a self-contradiction that can only end in error and confusion.


\flrightit{Posted on 2013-02-04 by Cologero}

\begin{center}* * *\end{center}

\begin{footnotesize}
\begin{sffamily}
\texttt{Jason-Adam on 2013-02-05 at 00:19 said:}

Brilliant article.

Do you think Evola's error may have arisen from the fascistic cult of the totalitarian leader like Mussolini or Hitler above which there is no power ?

\hfill

\texttt{Avery Morrow on 2013-02-05 at 01:52 said:}

Evola was opposed to those kinds of personality cults, but we have documented his strong interest in Nietzsche here.

\hfill

\texttt{Jason-Adam on 2013-02-05 at 13:02 said:}

Is there a specific reference in Nietzsche to justify one man having absolute power over both the sacred and secular realms ? Evola's writing on the King seems to me as a critique of the Holy Roman Empire's system where the Pope and Emperor always fought over where the boundaries between their powers ended and began. Guenon supported the Pope's claim to dominance but Evola would have preferred a warrior being in full command over the Pope, is that from Nietzsche ?

\hfill

\texttt{George on 2013-02-05 at 20:43 said:}

This helps clarify some of Jesus' sayings as well. The church (Christians) are wed to Jesus as His bride. Christians are commanded to show God and Jesus love, and to worship Him as Lord.

For a specific example, in Luke 22 where Jesus commands His followers to take up swords, He may have been delineated caste functions to some extent -- while He never took up a sword Himself (unmoved mover?). The passage in Luke tends to cause a great deal of confusion with modern Christians (who, for the most part, simply ignore anything they can't literally interpret to support leftism). In light of your article though, things make more sense.


\hfill

\end{sffamily}
\end{footnotesize}

